
\documentclass[10pt]{article}
\usepackage[utf8]{inputenc}
\usepackage[T1]{fontenc}
\usepackage{amsmath, amssymb, xcolor, geometry, enumitem, multicol, tcolorbox}

\definecolor{mainblue}{RGB}{0, 102, 204}
\geometry{a4paper, margin=1.2cm, top=1cm, bottom=3cm, footskip=1.2cm}

\begin{document}
\enlargethispage{2.5cm}

% Modern Header
\begin{tcolorbox}[colback=mainblue!5, colframe=mainblue, sharp corners, boxrule=0.5pt]
    \small \textbf{Unit:} Unit 0: Foundations of Algebra \hfill \textbf{Name:} \rule{3cm}{0.4pt} \\
    \small \textbf{Topic:} Natural Numbers & Integers \hfill \textbf{Date:} \rule{2cm}{0.4pt} \\
    \centering \textbf{\large Homework (Jalapeno)}
\end{tcolorbox}

\vspace{0.2cm}

% MCQ Section
\begin{multicols}{2}
\begin{enumerate}[label=\textbf{Q\arabic*.}, leftmargin=*, itemsep=6pt, parsep=0pt]

  \item \small The ancient Egyptians used $\mathbb{N}$ to represent which set of numbers?
  \begin{enumerate}[label=(\Alph*), leftmargin=0.6cm, nosep, itemsep=2pt]
    \item Whole numbers
    \item Integers
    \item Natural numbers
    \item Rational numbers
    \item Irrational numbers
  \end{enumerate}

  \item \small The number $5$ is classified as which type of number?
  \begin{enumerate}[label=(\Alph*), leftmargin=0.6cm, nosep, itemsep=2pt]
    \item Even and whole
    \item Odd and natural
    \item Even and integer
    \item Odd and integer
    \item Even and rational
  \end{enumerate}

  \item \small The ancient Babylonians used a sexagesimal system that included which of the following numbers?
  \begin{enumerate}[label=(\Alph*), leftmargin=0.6cm, nosep, itemsep=2pt]
    \item Only whole numbers
    \item Only integers
    \item Only natural numbers
    \item Whole numbers and fractions
    \item Integers and fractions
  \end{enumerate}

  \item \small The pattern $2, 5, 8, 11, 14, ...$ is an example of which type of sequence?
  \begin{enumerate}[label=(\Alph*), leftmargin=0.6cm, nosep, itemsep=2pt]
    \item Arithmetic sequence
    \item Geometric sequence
    \item Fibonacci sequence
    \item Harmonic sequence
    \item Alternating sequence
  \end{enumerate}

  \item \small The result of $(-3) + 5$ is which of the following?
  \begin{enumerate}[label=(\Alph*), leftmargin=0.6cm, nosep, itemsep=2pt]
    \item $-8$
    \item $-2$
    \item $2$
    \item $8$
    \item $12$
  \end{enumerate}

  \item \small The ancient Greeks used $\mathbb{Z}$ to represent which set of numbers?
  \begin{enumerate}[label=(\Alph*), leftmargin=0.6cm, nosep, itemsep=2pt]
    \item Whole numbers
    \item Integers
    \item Natural numbers
    \item Rational numbers
    \item Irrational numbers
  \end{enumerate}

  \item \small The number $-2$ is classified as which type of number?
  \begin{enumerate}[label=(\Alph*), leftmargin=0.6cm, nosep, itemsep=2pt]
    \item Even and whole
    \item Odd and natural
    \item Even and integer
    \item Odd and integer
    \item Even and rational
  \end{enumerate}

  \item \small The result of $(-2) \cdot (-5)$ is which of the following?
  \begin{enumerate}[label=(\Alph*), leftmargin=0.6cm, nosep, itemsep=2pt]
    \item $-10$
    \item $-4$
    \item $10$
    \item $4$
    \item $20$
  \end{enumerate}

\end{enumerate}
\end{multicols}

\vspace{-0.2cm}
\hrule
\vspace{0.2cm}
\textbf{\small Open Ended Questions (FRQ)}
\begin{enumerate}[label=\textbf{Q\arabic*.}, start=9, itemsep=5pt, topsep=0pt]

  \item \small The ancient Mayans used a vigesimal system to represent numbers. If a Mayan numeral represents the number $x$, and $x$ is a natural number, show that $x$ can be represented as $x = a_0 + a_1 \cdot 20 + a_2 \cdot 20^2 + ...$ where $a_i$ are digits in the Mayan system. \vspace{2cm}

  \item \small A historian discovers an ancient scroll with a pattern of numbers $1, 4, 7, 10, 13, 16, ...$. If the pattern continues, what is the next number in the sequence? Show your work and explain your reasoning. \vspace{2cm}

\end{enumerate}

\vfill

% Chef's Tips Footer
\begin{tcolorbox}[colback=blue!5, colframe=mainblue!50, title=\small \textbf{Chef's Tips}, fonttitle=\bfseries, bottom=1mm]
\begin{itemize}[noitemsep, topsep=2pt, leftmargin=0.5cm]
  \item \scriptsize Tip 1: Allow students to use calculators for checkwork only.\item \scriptsize Tip 2: Emphasize the importance of showing work for free-response questions.\item \scriptsize Tip 3: Encourage students to use historical context to help solve problems.
\end{itemize}
\end{tcolorbox}

\end{document}
  